\begin{examplebox}
	\textbf{\textcolor{CExample}{例1：利用三角形中位线}}

	\textbf{题目：}

	在空间四边形 $ABCD$ 中，$E,F$ 分别是 $AB,AD$ 的中点。求证：
	\[
		EF\parallel\text{平面 }BCD.
	\]

	\textbf{证明：}

	\begin{description}[
			style=nextline,
			leftmargin=2.8em,
			labelsep=0.5em,
			itemsep=0.45em
		]
		\item[\textbf{第一步（寻找面内参照线）：}]
		      连接 $BD$。

		      因为 $B,D$ 都在平面 $BCD$ 内，所以
		      \[
			      BD\subset\text{平面 }BCD.
		      \]

		\item[\textbf{第二步（证明线线平行）：}]
		      在 $\triangle ABD$ 中，$E,F$ 分别是 $AB,AD$ 的中点。

		      由三角形中位线定理得
		      \[
			      EF\parallel BD.
		      \]

		\item[\textbf{第三步（确认目标线不在平面内）：}]
		      因为 $ABCD$ 是空间四边形，所以点 $A$ 不在平面 $BCD$ 内。

		      点 $E$ 是线段 $AB$ 的中点。

		      假设 $E\in\text{平面 }BCD$，由于
		      \[
			      B\in\text{平面 }BCD,
		      \]
		      则直线 $BE$ 在平面 $BCD$ 内。

		      又因为 $A,B,E$ 共线，所以直线 $BE$ 就是直线 $AB$，从而
		      \[
			      A\in\text{平面 }BCD,
		      \]
		      产生矛盾。

		      因此
		      \[
			      E\notin\text{平面 }BCD,
		      \]
		      所以
		      \[
			      EF\not\subset\text{平面 }BCD.
		      \]

		\item[\textbf{第四步（套用判定定理）：}]
		      因为
		      \[
			      EF\not\subset\text{平面 }BCD,
		      \]
		      \[
			      BD\subset\text{平面 }BCD,
		      \]
		      且
		      \[
			      EF\parallel BD,
		      \]
		      所以由直线与平面平行的判定定理得
		      \[
			      EF\parallel\text{平面 }BCD.
		      \]
	\end{description}

	\medskip
	\textbf{\textcolor{CExample}{例2：利用平行四边形与中位线}}

	\textbf{题目：}

	在四棱锥 $P-ABCD$ 中，底面 $ABCD$ 为平行四边形，$E$ 为 $PC$ 的中点。求证：
	\[
		PA\parallel\text{平面 }BDE.
	\]

	\textbf{证明：}

	\begin{description}[
			style=nextline,
			leftmargin=2.8em,
			labelsep=0.5em,
			itemsep=0.45em
		]
		\item[\textbf{第一步（构造面内参照线）：}]
		      连接 $AC$，设
		      \[
			      AC\cap BD=O.
		      \]

		      因为四边形 $ABCD$ 是平行四边形，所以它的对角线互相平分。

		      因此点 $O$ 是 $AC$ 的中点。

		      连接 $EO$。

		      因为
		      \[
			      E\in\text{平面 }BDE,
			      \qquad
			      O\in BD\subset\text{平面 }BDE,
		      \]
		      所以
		      \[
			      EO\subset\text{平面 }BDE.
		      \]

		\item[\textbf{第二步（证明线线平行）：}]
		      在 $\triangle PCA$ 中，$E,O$ 分别是 $PC,CA$ 的中点。

		      由三角形中位线定理得
		      \[
			      EO\parallel PA.
		      \]

		\item[\textbf{第三步（确认目标线不在平面内）：}]
		      假设
		      \[
			      A\in\text{平面 }BDE.
		      \]

		      因为 $A,B,D$ 三点不共线，所以 $A,B,D$ 确定的平面就是底面 $ABCD$。

		      于是平面 $BDE$ 与底面 $ABCD$ 重合，从而
		      \[
			      E\in\text{平面 }ABCD.
		      \]

		      但 $E$ 是侧棱 $PC$ 的中点，且顶点 $P$ 不在底面 $ABCD$ 内，因此除点 $C$ 外，线段 $PC$ 上的点均不在底面内。

		      所以
		      \[
			      E\notin\text{平面 }ABCD,
		      \]
		      产生矛盾。

		      因此
		      \[
			      A\notin\text{平面 }BDE,
		      \]
		      从而
		      \[
			      PA\not\subset\text{平面 }BDE.
		      \]

		\item[\textbf{第四步（套用判定定理）：}]
		      因为
		      \[
			      PA\not\subset\text{平面 }BDE,
		      \]
		      \[
			      EO\subset\text{平面 }BDE,
		      \]
		      且
		      \[
			      PA\parallel EO,
		      \]
		      所以
		      \[
			      PA\parallel\text{平面 }BDE.
		      \]
	\end{description}

	\medskip
	\textbf{\textcolor{CExample}{例3：利用棱柱中的中位线}}

	\textbf{题目：}

	在直三棱柱 $ABC-A_1B_1C_1$ 中，$D$ 为 $BC$ 的中点。求证：
	\[
		A_1B\parallel\text{平面 }ADC_1.
	\]

	\textbf{证明：}

	\begin{description}[
			style=nextline,
			leftmargin=2.8em,
			labelsep=0.5em,
			itemsep=0.45em
		]
		\item[\textbf{第一步（构造面内参照线）：}]
		      连接 $A_1C$ 与 $AC_1$，设它们交于点 $O$。

		      因为侧面 $ACC_1A_1$ 是矩形，所以它的对角线互相平分。

		      因此点 $O$ 是 $A_1C$ 的中点，同时
		      \[
			      O\in AC_1.
		      \]

		      连接 $OD$。

		      因为
		      \[
			      O\in AC_1\subset\text{平面 }ADC_1,
		      \]
		      且
		      \[
			      D\in\text{平面 }ADC_1,
		      \]
		      所以
		      \[
			      OD\subset\text{平面 }ADC_1.
		      \]

		\item[\textbf{第二步（证明线线平行）：}]
		      在 $\triangle A_1BC$ 中，$O,D$ 分别是 $A_1C,BC$ 的中点。

		      由三角形中位线定理得
		      \[
			      OD\parallel A_1B.
		      \]

		\item[\textbf{第三步（确认目标线不在平面内）：}]
		      假设
		      \[
			      A_1\in\text{平面 }ADC_1.
		      \]

		      因为 $A_1,C_1$ 都在平面 $ADC_1$ 内，所以
		      \[
			      A_1C_1\subset\text{平面 }ADC_1.
		      \]

		      又因为棱柱的对应棱平行，所以
		      \[
			      AC\parallel A_1C_1.
		      \]

		      直线 $AC$ 经过平面 $ADC_1$ 内的点 $A$，且平行于平面内的直线 $A_1C_1$，因此
		      \[
			      AC\subset\text{平面 }ADC_1.
		      \]

		      从而
		      \[
			      C\in\text{平面 }ADC_1.
		      \]

		      于是平面 $ADC_1$ 与底面 $ABC$ 同时包含不共线的三点 $A,C,D$，所以两个平面重合。

		      这将导致
		      \[
			      C_1\in\text{平面 }ABC,
		      \]
		      与棱柱的上、下底面不重合矛盾。

		      因此
		      \[
			      A_1\notin\text{平面 }ADC_1,
		      \]
		      从而
		      \[
			      A_1B\not\subset\text{平面 }ADC_1.
		      \]

		\item[\textbf{第四步（套用判定定理）：}]
		      因为
		      \[
			      A_1B\not\subset\text{平面 }ADC_1,
		      \]
		      \[
			      OD\subset\text{平面 }ADC_1,
		      \]
		      且
		      \[
			      A_1B\parallel OD,
		      \]
		      所以
		      \[
			      A_1B\parallel\text{平面 }ADC_1.
		      \]
	\end{description}
\end{examplebox}